\section{Flag, Tanda Bilangan, dan Verifikasi Hasil}
\label{sec:flag-verifikasi}

Hasil aritmatika tidak hanya berupa angka pada register tujuan.
Prosesor juga mencatat keadaan hasil melalui \eng{flag}. Catatan
inilah yang dipakai program untuk menentukan apakah hasil masih sah,
apakah terjadi luapan, atau apakah dua nilai sama besar.

\subsection{CF untuk Tak Bertanda, OF untuk Bertanda}

\flag{CF} dan \flag{OF} sering membingungkan karena keduanya
menandai luapan. Perbedaannya terletak pada tafsiran bilangan.

\begin{itemize}
    \item \flag{CF} dipakai untuk bilangan tidak bertanda. Flag ini
          bernilai 1 bila hasil melewati batas atas 8 bit atau 16 bit.
    \item \flag{OF} dipakai untuk bilangan bertanda. Flag ini
          bernilai 1 bila hasil keluar dari jangkauan bilangan
          bertanda, misalnya dua bilangan positif menghasilkan
          bilangan negatif.
\end{itemize}

Contoh pada 8 bit:

\begin{lstlisting}[language=MASM, numbers=none]
MOV AL, 0FFh
ADD AL, 1           ; AL = 00h, CF = 1, OF = 0
MOV AL, 7Fh
ADD AL, 1           ; AL = 80h, CF = 0, OF = 1
MOV AL, 80h
ADD AL, 80h         ; AL = 00h, CF = 1, OF = 1
\end{lstlisting}

Kasus pertama adalah luapan tidak bertanda: 255 + 1 = 256, tidak muat
di 8 bit. Kasus kedua adalah luapan bertanda: 127 + 1 = 128, melewati
batas bilangan bertanda 8 bit. Kasus ketiga memicu keduanya.

\subsection{Flag Lain}

\begin{itemize}
    \item \flag{ZF} bernilai 1 bila hasilnya nol. Flag ini dipakai
          untuk memeriksa apakah dua operand sama besar setelah
          \asm{CMP}.
    \item \flag{SF} bernilai 1 bila bit tanda hasil menyala, yaitu
          hasil bertanda negatif.
    \item \flag{AF} menandai simpanan dari empat bit rendah. Flag ini
          dipakai pada operasi BCD.
    \item \flag{PF} menandai paritas byte rendah hasil operasi.
\end{itemize}

\asm{INC} dan \asm{DEC} memperbarui \flag{ZF}, \flag{SF},
\flag{OF}, \flag{AF}, dan \flag{PF}, tetapi tidak mengubah
\flag{CF}. Karena itu, keduanya tidak dapat menggantikan \asm{ADD}
atau \asm{SUB} pada aritmatika bertingkat.

\subsection{Verifikasi dengan Debugger}

Hasil operasi dapat diperiksa dengan \eng{single step} pada emu8086
atau DEBUG. Langkah yang dilakukan:

\begin{enumerate}
    \item jalankan program sampai instruksi aritmatika yang diamati;
    \item catat isi register tujuan sebelum dan sesudah instruksi;
    \item catat \flag{CF}, \flag{OF}, \flag{ZF}, dan \flag{SF};
    \item bandingkan hasilnya dengan perhitungan manual.
\end{enumerate}

Program \file{operasi-aritmatika.asm} menyimpan hasil setiap operasi
ke segmen data. Nilai tersebut dapat dibaca dengan perintah
\lstasm|D DS:0000|. Perkiraan hasilnya dirangkum pada
Tabel~\ref{tab:trace-aritmatika}.

\begin{table}[H]
    \centering
    \caption{Perkiraan hasil operasi pada
    \file{operasi-aritmatika.asm}.}
    \label{tab:trace-aritmatika}
    \small
    \begin{tabular}{@{}p{4.6cm}p{3.0cm}p{5.6cm}@{}}
        \toprule
        \textbf{Operasi} & \textbf{Hasil} &
        \textbf{Keterangan} \\
        \midrule
        \hex{0001FFFF} + \hex{00020001} &
        \hex{0004}:\hex{0000} &
        \asm{ADD} menghasilkan \flag{CF} = 1; \asm{ADC} membawanya \\
        \addlinespace
        \hex{00050000} $-$ \hex{00020001} &
        \hex{0002}:\hex{FFFF} &
        \asm{SUB} menghasilkan pinjaman; \asm{SBB} melanjutkannya \\
        \addlinespace
        $25 \times 4$ & \hex{0064} &
        \asm{MUL} 8 bit; hasil di \reg{AX} \\
        \addlinespace
        $1000 \div 7$ & \hex{008E}, sisa \hex{0006} &
        \asm{DIV} 8 bit; hasil bagi di \reg{AL}, sisa di \reg{AH} \\
        \bottomrule
    \end{tabular}
\end{table}

\begin{infobox}
\vspace{1.1em}
Bila hasil \asm{DIV} tidak sesuai, periksa lebih dahulu apakah
\reg{AH} (untuk pembagian 8 bit) atau \reg{DX} (untuk pembagian
16 bit) lebih kecil daripada pembagi. Kesalahan \eng{divide
overflow} membuat program berhenti sebelum hasil sempat disimpan.
\end{infobox}

\begin{peringatanbox}
\vspace{1.1em}
Jangan menyimpulkan hasil dari satu \eng{flag} saja. \flag{CF} = 1
belum tentu berarti hasil salah; itu hanya berarti hasil tidak
bertanda melewati kapasitas. Tafsirkan \eng{flag} sesuai jenis
bilangan yang sedang diolah.
\end{peringatanbox}
